\chapter{Обзор литературы}
\label{cha:Obzor}

   \section{Обзор научной литературы}
 
 В статьях \cite{TR-MFTI-khelvas2019-boxes, khel-EnT2019} решается задача трехмерной укладки коробок в паллету, относящаяся к классу NP-полных задач. Предложено несколько эвристических подходов к ее решению: на основе эвристик слоев и на основе генетического алгоритма. Предложен подход к оценке качества укладки на основе коэффициента перколяции и коэффициента устойчивости.
Показана слабая зависимость качества укладки от точности задания размеров коробок.

В статье  \cite{karabulut} представлен метод DBLF, который используется в гибридном генетическом алгоритме. Также для сравнения используются несколько эврестических методов. Алгоритм состоит из двух этапов: кодирования на основе порядка OBE для представления решений в виде хромосом и декодирования методом DBLF. Используются четыре функции приспособленности: максимальная высота контейнера, объем неиспользованного пространства, отношение упакованных и пустых объемов, комбинированная функция, учитывающая 70\% максимальной высоты и 30\% соотношения объемов.
Популяция обновляется с использованием рулеточного отбора, мутации осуществляются с вероятностью 0,05\%, а кроссовер с вероятностью 20\%.В результате проведённых экспериментов было установлено, что функция приспособленности, основанная на максимальной высоте, продемонстрировала наилучшие результаты.

В статье \cite{korelin} представлен модифицированный подход к классическому методу генетического алгоритма. Для селекции авторы используют два метода: турнирную селекцию и метод  рулетки. В качестве метода скрещивания применяются модификации одноточечного и двухточечного кроссинговера, при этом вторая часть потомка дополняется недостающими номерами блоков, которые заимствуются из второго родителя в порядке их появления. Мутация также подверглась изменению: вместо обычного изменения гена на противоположное значение, происходит его перемещение в начало приоритетного списка.

В работе \cite{TimSort_teor} исследуется оптимистическая сортировка,~— предшественник TimSort,~— позволяющая быстро упорядочивать заранее частично отсортированные наборы данных. Такая предварительная сортировка часто используется в задачах упаковки для ускорения построения начальных раскладок.

В статье \cite{Wei2011} предложена эвристика Skyline для двумерной strip-packing-задачи. Метод описывает  свободное пространство профилем «линии неба» и размещает объекты в первую подходящую впадину, обеспечивая $O(n\log n)$ сложность и хорошие результаты по плотности заполнения; идея впоследствии была обобщена на 3-D-упаковку.

Авторы \cite{Zhang2024} объединили генетический алгоритм с генеративной состязательной сетью (GAN), где GAN создаёт перспективные хромосомы, а GA эволюционирует их. Подход уменьшил невостребованный объём на 2–3 п.п. и ускорил сходимость по сравнению с чистым GA.

В работе \cite{Calzavara2021} сформулирована MIP-модель паллетизации с практическими ограничениями (выступы, устойчивость) и разработан гибрид GRASP + локальный поиск. На реальных наборах высота паллет снижена на 6\% по сравнению с промышленными правилами укладки.

Статья \cite{Queiroz2012} посвящена 3-D guillotine-задачам (knapsack, cutting-stock, strip-packing). Для них представлены динамические и колонковые алгоритмы.

Авторы \cite{Ananno2024} решают многоконтейнерную 3-D-BPP с реальными ограничениями, применяя двухэтапный multi-heuristic (зонная декомпозиция + шаблонный поиск). На производственных данных удалось сократить число требуемых контейнеров в среднем на один.

В \cite{Alonso2016} одновременно оптимизируется построение паллет и погрузка автомобилей в междеповской логистике; получена экономия транспорта 4–7 \% по сравнению с последовательным решением задач.

Работа \cite{Fang2023} внедряет глубокое обучение (Actor–Critic) для двумерного strip-packing; сеть обучается без учителя и оптимизирует длину ленты, превосходя классические эвристики на 3 п.п. — подход релевантен развитию RL-методов в 3-D.

Обзор \cite{Ali2022} систематизирует offline- и online-алгоритмы трёхмерной упаковки, классифицируя ограничения и метрики задержки, а также предоставляет открытый набор тестов для online-BPP.

В \cite{Zhao2021} предлагается RL-политика для online 3-D-упаковки с гарантией устойчивости на основе «stacking-tree», что дало выигрыш до 15 \% по числу контейнеров относительно rule-based baseline.

Статья \cite{Amossen2010} вводит guillotine-ограничения в многомерное упаковывание и предлагает CP-модель с конструктивной эвристикой, показывая рост расхода объёма всего на 0,5 \% при полной технологичности разрезов.

Авторы \cite{Yang2023} решают задачу выбора «открытых» размеров коробок для e-commerce. MIP-модель 3D-MODRPP позволяет сократить расход картона на 8 \% и «воздушный» объём — на 12 \%.

В работе \cite{Faroe2003} применён Guided Local Search; штрафы за перегруженные измерения постепенно уменьшаются, что позволяет приближаться к нижней границе, сокращая число контейнеров на 0,4 \% за итерацию. 

Алгоритм Maximal-Space \cite{Parreno2008} описывает свободный объём множеством перекрывающихся «max-spaces» и обновляет их после каждой укладки; метод улучшил лучшие известные решения на 25 из 30 задач Bischoff–Ratcliff.

Методологическая работа \cite{Bischoff1995} раскрывает влияние практических ограничений (несовместимость грузов, очаговые нагрузки) на выбор модели и эвристики, оставаясь актуальным справочником.

Типология \cite{Wascher2007} расширяет классификацию Dyckhoff, выделяя четыре уровня описания задач cutting \& packing и став «золотым стандартом» при формулировке новых постановок.

Обзор \cite{Bortfeldt2013} концентрируется на ограничениях контейнерной загрузки и закрывает период 1995–2013 гг., обеспечивая полный исторический контекст.

Статья \cite{Martello2000} сочетает branch-and-bound с эвристикой L2, достигая оптимума для задач с до 100 коробок; её инстансы используются как бенчмарк и по сей день.

Практическое руководство \cite{Jylanki2010} сравнивает более 20 двумерных эвристик и предлагает C++-библиотеку RectangleBinPack, широко применяемую в игровой и GUI-индустрии.


\AKHworries{Обзор Светлов }

Работа \cite{Lodi2002} предлагает эвристику трёхмерной упаковки паллеты, основанную на послойном построении решения с предварительным выбором высоты слоя и двумерным размещением коробок внутри слоя. Подход обеспечивает эффективное вертикальное заполнение контейнера и используется как базовый конструктивный метод в метаэвристиках.

В статье \cite{Crainic2007} введено правило Extreme Points для трёхмерной упаковки, определяющее кандидатные позиции размещения на границе уже упакованных объектов. Метод сохраняет слоистую структуру построения решения и демонстрирует более высокую плотность заполнения по сравнению с классическими shelf- и corner-point-подходами.

Обзор \cite{Zhao2014} систематизирует алгоритмы трёхмерной загрузки контейнеров и классифицирует placement-эвристики, включая layer-building и extreme-point-методы. Сделан вывод о том, что методы с предвыбором высоты слоя используются в большинстве практических решений.


\AKHworries{Обзор Юхневич }

В работе \cite{Pisinger2005} представлен метод Large Neighborhood Search (LNS) --- метаэвристика, которая итеративно разрушает и восстанавливает часть текущего решения. LNS показал отличные результаты на задачах vehicle routing и bin packing, превосходя классические локальные поиски за счёт возможности выхода из локальных оптимумов.

Статья \cite{Ropke2006} развивает идею Adaptive LNS (ALNS), где веса операторов разрушения и восстановления адаптивно подстраиваются в процессе поиска на основе их эффективности. На задачах pickup-and-delivery ALNS улучшил лучшие известные решения на 40\% инстансов.

В работе \cite{Shaw2002} впервые предложен constraint-based подход к LNS, где выбор разрушаемых переменных основан на их взаимосвязи с <<плохими>> частями решения. Идея <<связанности>> (relatedness) стала основой многих современных операторов разрушения.

Обзор \cite{Pisinger2010} систематизирует развитие LNS, выделяя ключевые компоненты: выбор neighbourhood size, стратегии acceptance (simulated annealing, threshold accepting, record-to-record travel), и hybrid-методы с exact solvers.

В статье \cite{Ekici2023} LNS применён к задаче bin packing с дополнительными ограничениями (конфликты между предметами). Авторы показали, что LNS с адаптивными операторами превосходит другие метаэвристики на задачах средней и большой размерности.


\AKHworries{Обзор Галкин  }


\section{Патентный обзор}
В патенте \cite{USPatent8423178B2} предложен метод формирования сборной палетты из коробок находящихся в множестве однородных паллет.

В патенте \cite{US11613434B2} описан способ послойного формирования палетты. 